\section*{LỜI NÓI ĐẦU}
\phantomsection
\addcontentsline{toc}{section}{\numberline{}LỜI NÓI ĐẦU}

Giới thiệu tên đề tài, chú ý \textbf{"bôi đậm tên đề tài"}, đồng thời nêu bối cảnh và lý do lựa chọn đề tài, nêu bật tầm quan trọng hoặc tính cần thiết của vấn đề trong lĩnh vực nghiên cứu. Trình bày ngắn gọn mục tiêu thực hiện, phạm vi đề tài và động lực thực hiện (nếu phù hợp). Nội dung cần được trình bày trang trọng, chân thành và súc tích; tránh đi quá sâu vào các nội dung kỹ thuật hoặc trình bày chi tiết kết quả, vì những nội dung này sẽ được đề cập trong các phần tiếp theo của báo cáo. Lưu ý: Độ dài mục lời nói đầu không quá 1 trang.

Có thể gửi lời cảm ơn đến giảng viên hướng dẫn, các cá nhân hoặc tổ chức đã hỗ trợ trong quá trình thực hiện.

...

Chúng em đã rất cố gắng để hoàn thành tốt nhất đề tài của mình. Tuy nhiên, do vốn kiến thức còn chưa nhiều, cũng như nhiều yếu tố khách quan tác động mà đề tài của nhóm còn nhiều thiếu sót. Chúng em rất mong thầy/cô đóng góp ý kiến, phê bình và hướng dẫn để các đề tài sau được hoàn thiện hơn.

Chúng em xin gửi lời cảm ơn chân thành nhất tới ... đã hướng dẫn tận tình, chi tiết để chúng em có thể hoàn thành tốt đề tài này. Qua đây, chúng em không chỉ có thêm nhiều kiến thức về chuyên môn mà còn nâng cao hơn được kỹ năng làm việc nhóm, trau dồi thêm ngoại ngữ, biết cách phân tích và tư duy, cùng nhiều kỹ năng mềm khác.

\begin{flushright}
    \noindent\textit{Chúng em xin chân thành cảm ơn!}
\end{flushright}

...

\newpage
\cleardoublepage